\begin{center}\bfseries e.\quad\textls[100]{Kreisteilung}\footnote{Vergl. die einschlägigen Abschnitte der weiter unten folgenden Aufsätze »Über Gauss’ algebraische Arbeiten« von K. \textsc{Hensel} und »Gauss als Geometer« von P. \textsc{Stäckel}.}.\end{center}
\begin{center}\bfseries 13.\end{center}
\noindent\hspace*{\parindent}Der \textls[100]{Fortgang zum siebenten Abschnitte} der D. A. führt uns nun zu Ergebnissen, welche wohl unter \textsc{Gauss}’ arithmetischen Entdeckungen das vielseitigste Interesse darbieten. Eigentlich nur eine ganz einfache Anwendung eines elementaren arithmetischen Satzes, sind sie gleichwohl zur Quelle der schönsten zahlentheoretischen Erkenntnisse geworden, ja schliesslich zur Grundlage der stolzesten Gebilde, welche die neuere Zahlentheorie in ihrer Lehre von den Zahlkörpern und deren Idealen geschaffen hat. Andererseits aber hat eben jene Anwendung die engste Verbindung herbeigeführt zwischen der höheren Arithmetik auf der einen, und der Theorie der algebraischen Gleichungen auf der anderen Seite, und wieder durch die Natur der in Frage kommenden Gleichungen jene Wissenschaft auch als Ursprung tiefliegender Wahrheiten der Geometrie erwiesen, eines Gebietes, das an sich gänzlich ihr
\newpage\SourceLeaf{40}\ArithmeticHead{33}{R}{KREISTEILUNG.}\HistoricalFoot{X\textsubscript{2}.}{5}
\noindent abseits zu liegen schien. Hatte die Lösung der Aufgabe, die Kreisperipherie in eine Anzahl gleicher Teile zu teilen, völlig auf dem Standpunkte beharrt, den sie vor 2000 Jahren schon erreicht und auf dem sie scheinbar sich erschöpft hatte, so lehrte jene arithmetische Anwendung neben den schon bekannten noch mancherlei Teiler kennen, für welche die Aufgabe ebenfalls mit Zirkel und Lineal ausführbar ist, insbesondere die Konstruktion des regelmässigen Siebzehnecks. Wie ein neuer glänzender Stern am Firmamente musste diese Entdeckung Aufsehen erregen, und \textsc{Gauss} selbst schätzte sie gebührend hoch ein, indem er damit sein \textit{Tagebuch} begann und von seinem am 30. März 1796 (Nr. 1 T.) gefundenen Forschungsresultate in Nr. 66 des Intelligenzblatts der allgemeinen Litteraturzeitung vom 1. Juni 1796 (W. X 1, S. 3) unter dem Datum des 18. April schon vor der Veröffentlichung seiner erst teilweise fertigen Theorie Kenntnis gab.

Wie \textsc{Gauss} zu seiner Theorie geführt sein mag? Kaum wohl von Seiten der Geometrie, der sein Interesse, will man \textsc{Sartorius v. Waltershausen} glauben\footnote{\textit{Gauss zum Gedächtniss}, Leipzig 1856, S. 80, 81.}, in seiner Jugend weniger zugewendet war, wahrscheinlicher durch algebraische Studien, welchen er ohne Zweifel neben seinen arithmetischen schon frühzeitig obgelegen hat und aus denen dann seine Doktorarbeit (1799, W. III, S. 1) hervorgegangen ist. Welche Gleichung höheren Grades konnte ihm näher liegen, als die Gleichung \(x^p=1\), deren Wurzeln dann das Problem der Kreisteilung von selber in den Gesichtskreis zogen! Wie bedeutsam erscheint dabei, dass \textsc{Gauss} diese einfachere Gleichung jenen andern, durch welche die trigonometrischen Funktionen von \(\frac{2\pi}{p}\) selbst bestimmt werden, vorzog, obwohl ihre Wurzeln zu den imaginären Grössen zählen, deren Bürgerrecht in der Wissenschaft damals noch so umstritten war, dass er sie in seiner Doktorarbeit geflissentlich vermied.

Bedenkt man aber, dass in Kap. 6 der A. R. die Theorie der \textls[100]{Gleichung} \(x^p=1\) erst derjenigen der \textls[100]{Kongruenz} \(x^p\equiv1\) nachfolgt und dass die Methode zur Auflösung dieser Kongruenz das genaue Vorbild für die Auflösung der Kreisteilungsgleichung ist, so darf man wohl auf den rein \textls[100]{arithmetischen} Ursprung von \textsc{Gauss}’ Beschäftigung mit der letztern schliessen. Ist doch das Prinzip seiner Methode von dieser Art. Die Wurzeln jener Gleichung oder
\newpage\SourceLeaf{41}\ArithmeticHead{34}{L}{BACHMANN, ÜBER GAUSS’ ZAHLENTHEORETISCHE ARBEITEN.}
\noindent genauer der Gleichung
\[
X=\frac{x^p-1}{x-1}=0,
\]
die, wenn wir
\[
r=\cos\frac{2\pi}{p}+\sqrt{-1}\sin\frac{2\pi}{p}
\]
setzen, durch die Potenzen \(r,r^2,r^3,\ldots,r^{p-1}\) dargestellt sind, zeigen, wenn \(p\) als eine ungerade Primzahl vorausgesetzt wird, auf Grund des arithmetischen Satzes, dass die Potenzen \(1,g,g^2,\ldots g^{p-2}\) einer primitiven Wurzel \(g\ (\mathrm{mod.}\ p)\) den Resten \(1,2,3,\ldots,p-1\), von der Ordnung abgesehen, kongruent sind, das eigentümliche Verhalten, dass bei der zyklischen Anordnung \(r,r^g,r^{g^2},\ldots,r^{g^{p-2}}\) jede die gleiche rationale Funktion (nämlich die \(g^{\mathrm{te}}\) Potenz) der vorhergehenden ist. \textls[100]{Mit dieser Einsicht war der schöpferische Gedanke\linebreak{} gewonnen, aus welchem die ganze \textsc{Gauss}sche Theorie der Kreisteilungsgleichung entsprang.}

Die Verallgemeinerung der \textsc{Gauss}schen Betrachtung, welche wir \textsc{Abel} verdanken\footnote{Siehe \textsc{Crelles} Journal für Mathematik \textbf{4} (1829), S. 131, Oeuvres, nouvelle édition I, S. 478.}, liess später die algebraische Auflösbarkeit aller Gleichungen erkennen, deren Wurzeln analoge Eigenschaft zeigen, und eröffnete den Weg zu den bahnbrechenden Ergebnissen von \textsc{Galois} und \textsc{Kronecker} über die algebraische Auflösung der Gleichungen, deren wir uns jetzt erfreuen. Übrigens war sich \textsc{Gauss} selbst schon zur Zeit der D. A. bewusst, dass seine Methode zur Auflösung der Kreisteilungsgleichung viel weiter reiche. Im art. 335 weist er in dieser Hinsicht auf die Funktionen hin, die aus dem Integrale
\[
\int\frac{dx}{\sqrt{1-x^4}}
\]
entspringen und für die Lemniskate gleiche Bedeutung haben, wie die trigonometrischen für den Kreis, und er hat nach den Angaben seines Tagebuchs schon am 19. März 1797 (Nr. 60 T.) die Gleichung untersucht, die zur Teilung der Lemniskate dient, und am 21. desselben Monats (Nr. 62 T.) die Tatsache festgestellt, dass die Fünfteilung auch für diese Kurve geometrisch mittels Zirkel und Lineal ausführbar ist.
